\documentclass{article}

\usepackage[spanish]{babel}
\usepackage[utf8]{inputenc}
\usepackage{fancyhdr}
\usepackage{extramarks}
\usepackage{amsmath}
\usepackage{amsthm}
\usepackage{amsfonts}
\usepackage{tikz}
\usepackage{algorithm}
\usepackage{algpseudocode}
\usepackage{hyperref}
\usepackage{graphicx}
\usepackage{float}
\usepackage{natbib}
\usepackage{subcaption}

\usetikzlibrary{automata,positioning}

%
% Basic Document Settings
%

\topmargin=-0.45in
\evensidemargin=0in
\oddsidemargin=0in
\textwidth=6.5in
\textheight=9.0in
\headsep=0.25in

\linespread{1.1}

\pagestyle{fancy}
\lhead{\hmwkAuthorName}

\setlength{\headheight}{35pt}

\chead{
	\parbox{0.8\textwidth}{
		\centering
		\hmwkClass\ (\hmwkClassInstructor)\\
		\hmwkTitle
}}

\rhead{\firstxmark}
\lfoot{\lastxmark}
\cfoot{\thepage}

\renewcommand\headrulewidth{0.4pt}
\renewcommand\footrulewidth{0.4pt}

\setlength\parindent{0pt}

%
% Create Problem Sections
%

\newcommand{\enterProblemHeader}[2]{
	\nobreak\extramarks{}{Problema {#1} continúa en la sig pág\ldots}\nobreak{}
	\nobreak\extramarks{Problema \arabic{#1} (continúa)}{Problema \arabic{#1} continúa en la sig pág\ldots}\nobreak{}
}

\newcommand{\exitProblemHeader}[2]{
	\nobreak\extramarks{Problema \arabic{#1}}{Problema \arabic{#1} continúa en la sig pág\ldots}\nobreak{}
	\stepcounter{#1}
	\nobreak\extramarks{Problema \arabic{#1}}{}\nobreak{}
}

\setcounter{secnumdepth}{0}
\newcounter{partCounter}
\newcounter{homeworkProblemCounter}
\setcounter{homeworkProblemCounter}{8}
\nobreak\extramarks{Problema \arabic{homeworkProblemCounter}}{}\nobreak{}

%
% Homework Problem Environment
%
% This environment takes an optional argument. When given, it will adjust the
% problem counter. This is useful for when the problems given for your
% assignment aren't sequential. See the last 3 problems of this template for an
% example.
%
\newenvironment{homeworkProblem}[2][-1]{
	\ifnum#1>0
	\setcounter{homeworkProblemCounter}{#1}
	\fi
	\def\currentProblemTitle{#2}
	\section{Problema \arabic{homeworkProblemCounter}: #2}
	\setcounter{partCounter}{1}
	\enterProblemHeader{homeworkProblemCounter}{#2}
}{
	\exitProblemHeader{homeworkProblemCounter}{\currentProblemTitle}
}
%
% Homework Details
%   - Title
%   - Due date
%   - Class
%   - Section/Time
%   - Instructor
%   - Author
%

\newcommand{\hmwkTitle}{La Matriz de Denavit–Hartenberg\ \#8}
\newcommand{\hmwkDueDate}{21 de mayo de 2026}
\newcommand{\hmwkClass}{Planificación de movimientos de robots}
\newcommand{\hmwkClassInstructor}{Rafael Murrieta Cid}
\newcommand{\hmwkAuthorName}{\textbf{Andrés Alemán}}

%
% Title Page
%

\title{
	\vspace{2in}
	\textmd{\textbf{\hmwkClass:\ \hmwkTitle}}\\
	\normalsize\vspace{0.1in}\small{Due\ on\ \hmwkDueDate}\\
	\vspace{0.1in}\large{\textit{\hmwkClassInstructor}}
	\vspace{3in}
}

\author{\hmwkAuthorName}
\date{}

\renewcommand{\part}[1]{\textbf{\large Parte \Alph{partCounter}}\stepcounter{partCounter}\\}

%
% Various Helper Commands
%

% Useful for algorithms
\newcommand{\alg}[1]{\textsc{\bfseries \footnotesize #1}}

% For derivatives
\newcommand{\deriv}[1]{\frac{\mathrm{d}}{\mathrm{d}x} (#1)}

% For partial derivatives
\newcommand{\pderiv}[2]{\frac{\partial}{\partial #1} (#2)}

% Integral dx
\newcommand{\dx}{\mathrm{d}x}

% Alias for the Solution section header
\newcommand{\solution}{\textbf{\large Solución}}

% Probability commands: Expectation, Variance, Covariance, Bias
\newcommand{\E}{\mathrm{E}}
\newcommand{\Var}{\mathrm{Var}}
\newcommand{\Cov}{\mathrm{Cov}}
\newcommand{\Bias}{\mathrm{Bias}}

\begin{document}
	
%	\maketitle
	
%	\pagebreak
	
	\begin{homeworkProblem}{La Matriz de Denavit–Hartenberg}
	
	\section{Introducción}
	
	El problema fundamental de la cinemática de manipuladores consiste en describir la posición y orientación del efector final en función de las variables articulares, sin considerar las fuerzas que originan el movimiento \cite{murrieta2026}. Para simplificar este análisis, se utiliza comúnmente la convención de Denavit-Hartenberg (DH), la cual permite reducir los seis parámetros necesarios para describir un objeto en el espacio 3D a solo cuatro parámetros por eslabón: longitud de eslabón ($a_i$), torsión de eslabón ($\alpha_i$), desfasamiento de eslabón ($d_i$) y ángulo de articulación ($\theta_i$) \cite{spong2020}.
	
	Esta representación sistemática se basa en la asignación de marcos de referencia que deben cumplir con dos restricciones geométricas fundamentales: 
	\begin{enumerate}
		\item \textbf{DH1:} El eje $x_i$ es perpendicular al eje $z_{i-1}$.
		\item \textbf{DH2:} El eje $x_i$ intersecta al eje $z_{i-1}$.
	\end{enumerate}
	
	Bajo estas premisas, la transformación entre eslabones se define como el producto de cuatro transformaciones básicas: $A_i = R_{z,\theta_i} Trans_{z,d_i} Trans_{x,a_i} R_{x,\alpha_i}$ \cite{spong2020}. El presente documento tiene como objetivo completar la validación matemática de todos los elementos de la matriz resultante $A_i$, demostrando que su forma trigonométrica es una consecuencia directa de las restricciones de la convención DH y de las propiedades de ortogonalidad inherentes a las matrices de rotación del grupo $SO(3)$ \cite{murrieta2026, spong2020}.
					
	\subsection{La Matriz de Transformación resultante $A_i$}
	Como se mencionó anteriormente, la transformación total entre marcos de referencia sucesivos se obtiene mediante el producto de las cuatro transformaciones básicas definidas por los parámetros DH \cite{spong2020}. La matriz resultante, que es el objetivo de validación de este trabajo, tiene la siguiente forma:
	
	\begin{equation}
		A_i = \begin{pmatrix} 
			\cos \theta_i & -\sin \theta_i \cos \alpha_i & \sin \theta_i \sin \alpha_i & a_i \cos \theta_i \\ 
			\sin \theta_i & \cos \theta_i \cos \alpha_i & -\cos \theta_i \sin \alpha_i & a_i \sin \theta_i \\ 
			0 & \sin \alpha_i & \cos \alpha_i & d_i \\ 
			0 & 0 & 0 & 1 
		\end{pmatrix}
	\end{equation}
	
	Para facilitar la demostración de los elementos individuales, es útil representar la submatriz de rotación $R_i^{i-1}$ mediante sus componentes genéricos $r_{m,n}$ \cite{murrieta2026}. De este modo, la matriz de transformación homogénea se expresa como:
	
	\begin{equation}
		A_i = \begin{pmatrix} 
			r_{1,1} & r_{1,2} & r_{1,3} & a_i \cos \theta_i \\ 
			r_{2,1} & r_{2,2} & r_{2,3} & a_i \sin \theta_i \\ 
			r_{3,1} & r_{3,2} & r_{3,3} & d_i \\ 
			0 & 0 & 0 & 1 
		\end{pmatrix}
	\end{equation}
					
	\section{Demostración Parcial Realizada en Clase}
	
	La validación parcial presentada en la sesión de clase se centró en demostrar que los parámetros DH restringen la forma de la matriz de transformación $A_i$ de tal manera que se cumplen las condiciones geométricas del sistema \cite{murrieta2026}. A continuación, se detallan los elementos ya verificados:
	
	\subsection{Validación de $r_{3,1}$ mediante la restricción DH1}
	La primera restricción (DH1) establece que el eje $x_i$ debe ser perpendicular al eje $z_{i-1}$. En términos de álgebra lineal, esto implica que el producto punto entre sus vectores unitarios debe ser cero \cite{spong2020}. Al expresar esta relación en el marco de referencia $i-1$, tenemos:
	\begin{equation}
		x_i^{i-1} \cdot z_{i-1}^{i-1} = 
		\begin{pmatrix} r_{1,1} \\ r_{2,1} \\ r_{3,1} \end{pmatrix} \cdot 
		\begin{pmatrix} 0 \\ 0 \\ 1 \end{pmatrix} = r_{3,1} = 0
	\end{equation}
	Este resultado confirma que el elemento de la tercera fila y primera columna de la submatriz de rotación debe ser idénticamente nulo \cite{murrieta2026}.
	
	\subsection{Determinación de elementos mediante magnitud unitaria}
	Dado que cada columna de una matriz de rotación pertenece al grupo $SO(3)$, su magnitud debe ser unitaria. Sabiendo que $r_{3,1}=0$, la primera columna debe cumplir $r_{1,1}^2 + r_{2,1}^2 = 1$, lo que se satisface con la asignación trigonométrica:
	\begin{equation}
		(r_{1,1}, r_{2,1}) = (\cos \theta_i, \sin \theta_i)
	\end{equation}
	De manera análoga, para que la tercera fila mantenga su naturaleza unitaria considerando $r_{3,1}=0$, se deduce que $r_{3,2}^2 + r_{3,3}^2 = 1$, justificando la forma:
	\begin{equation}
		(r_{3,2}, r_{3,3}) = (\sin \alpha_i, \cos \alpha_i)
	\end{equation}
	
	\subsection{Validación del vector de traslación mediante DH2}
	La restricción DH2 indica que el eje $x_i$ intersecta al eje $z_{i-1}$. Esto permite modelar el desplazamiento desde el origen $O_{i-1}$ hasta $O_i$ como una traslación de distancia $d_i$ sobre $z_{i-1}$ seguida de una traslación de distancia $a_i$ sobre $x_i$ \cite{murrieta2026}. Expresado matemáticamente:
	\begin{equation}
		O_i^{i-1} = O_{i-1}^{i-1} + d_i z_{i-1}^{i-1} + a_i x_i^{i-1}
	\end{equation}
	Sustituyendo los valores conocidos, se obtiene la cuarta columna de la matriz $A_i$:
	\begin{equation}
		O_i^{i-1} = \begin{pmatrix} 0 \\ 0 \\ 0 \end{pmatrix} + d_i \begin{pmatrix} 0 \\ 0 \\ 1 \end{pmatrix} + a_i \begin{pmatrix} \cos \theta_i \\ \sin \theta_i \\ 0 \end{pmatrix} = \begin{pmatrix} a_i \cos \theta_i \\ a_i \sin \theta_i \\ d_i \end{pmatrix}
	\end{equation}
					
	\section{Desarrollo: Demostración de los elementos restantes}
	
	Para completar la validación de la matriz $A_i$, es necesario determinar los cuatro elementos centrales de la submatriz de rotación: $r_{1,2}, r_{2,2}, r_{1,3}$ y $r_{2,3}$. Esta deducción se fundamenta en que toda matriz de rotación $R \in SO(3)$ debe satisfacer las restricciones de magnitud y perpendicularidad entre sus columnas \cite{spong2020}.
	
	\subsection{Propiedades fundamentales de las matrices de rotación}
	De acuerdo con \cite{spong2020}, los elementos $r_{i,j}$ de la matriz están sujetos a dos restricciones matemáticas clave:
	\begin{enumerate}
		\item \textbf{Magnitud unitaria (Ecuación 2.25):} La suma de los cuadrados de los elementos de cada columna debe ser igual a la unidad:
		\begin{equation}
			\sum_{i=1}^{3} r_{ij}^2 = 1, \quad j \in \{1, 2, 3\} \label{eq:spong225}
		\end{equation}
		\item \textbf{Ortogonalidad (Ecuación 2.26):} El producto punto entre cualquier par de columnas distintas debe ser cero:
		\begin{equation}
			r_{1i}r_{1j} + r_{2i}r_{2j} + r_{3i}r_{3j} = 0, \quad i \neq j \label{eq:spong226}
		\end{equation}
	\end{enumerate}
	
	\subsection{Deducción de la segunda columna ($r_{1,2}$ y $r_{2,2}$)}
	Para hallar los elementos de la segunda columna, aplicamos primero la restricción de ortogonalidad \eqref{eq:spong226} entre la Columna 1 y la Columna 2:
	\begin{equation}
		(r_{1,1})(r_{1,2}) + (r_{2,1})(r_{2,2}) + (r_{3,1})(r_{3,2}) = 0
	\end{equation}
	Sustituyendo los valores conocidos ($\cos \theta_i, \sin \theta_i, 0$ y $\sin \alpha_i$):
	\begin{equation}
		\cos \theta_i \cdot r_{1,2} + \sin \theta_i \cdot r_{2,2} + 0 \cdot \sin \alpha_i = 0 \implies \cos \theta_i \cdot r_{1,2} = -\sin \theta_i \cdot r_{2,2}
	\end{equation}
	Esta relación sugiere una solución de la forma $r_{1,2} = -\sin \theta_i \cdot C$ y $r_{2,2} = \cos \theta_i \cdot C$. Para determinar el factor $C$, empleamos la condición de magnitud unitaria \eqref{eq:spong225}:
	\begin{equation}
		(-\sin \theta_i \cdot C)^2 + (\cos \theta_i \cdot C)^2 + (\sin \alpha_i)^2 = 1
	\end{equation}
	\begin{equation}
		C^2(\sin^2 \theta_i + \cos^2 \theta_i) + \sin^2 \alpha_i = 1 \implies C^2 = 1 - \sin^2 \alpha_i = \cos^2 \alpha_i
	\end{equation}
	Tomando la solución para un sistema de mano derecha, obtenemos $C = \cos \alpha_i$, lo que valida que:
	\begin{equation}
		r_{1,2} = -\sin \theta_i \cos \alpha_i, \quad r_{2,2} = \cos \theta_i \cos \alpha_i
	\end{equation}
	
	\subsection{Deducción de la tercera columna ($r_{1,3}$ y $r_{2,3}$)}
	Siguiendo un procedimiento análogo para la tercera columna mediante el producto punto entre la Columna 1 y la Columna 3:
	\begin{equation}
		\cos \theta_i \cdot r_{1,3} + \sin \theta_i \cdot r_{2,3} + 0 \cdot \cos \alpha_i = 0
	\end{equation}
	Bajo la condición de magnitud unitaria para la Columna 3 y considerando que $r_{3,3} = \cos \alpha_i$, se llega a la solución:
	\begin{equation}
		r_{1,3} = \sin \theta_i \sin \alpha_i, \quad r_{2,3} = -\cos \theta_i \sin \alpha_i
	\end{equation}
	Este resultado completa la verificación de todos los elementos de la matriz de rotación parcial $R_i^{i-1}$ presentada en las fuentes \cite{murrieta2026, spong2020}.
						
	\bibliographystyle{plainnat}
	\bibliography{biblio}
	
	\end{homeworkProblem}
	
\end{document}
